\section{Design simulační studie}
Pro ověření vlastností Woodruffovy metody byla provedena Monte Carlo simulace v prostředí R. Cílem simulace je zjistit, jak přesně metoda odhaduje rozptyl a zda sestrojené intervaly spolehlivosti skutečně pokrývají populační parametr s deklarovanou pravděpodobností 95 \%.

\subsection{Populace a generování dat}
Jako model pro generování dat bylo pro svou relevanci v ekonomických aplikacích (příjmová rozdělení) zvoleno **Log-normální rozdělení** $LN(\mu, \sigma)$.
Byl fixován parametr polohy $\mu = 0$ (v logaritmické škále).
Parametr tvaru $\sigma$ nabýval hodnot $\{0.5, 1.0, 1.5\}$.
\begin{itemize}
    \item $\sigma=0.5$: Mírná šikmost, blízké normálnímu rozdělení.
    \item $\sigma=1.5$: Extrémní šikmost a těžké chvosty.
\end{itemize}

\subsection{Sledované parametry}
Zajímáme se o odhad dvou kvantilů:
\begin{itemize}
    \item \textbf{Medián ($p=0.5$)}: Robustní míra polohy.
    \item \textbf{90. percentil ($p=0.9$)}: Chvostový kvantil, kde je hustota pravděpodobnosti nízká.
\end{itemize}
Simulace byla provedena pro rozsahy výběru $n \in \{200, 500, 1000, 2000\}$. Počet replikací byl nastaven na $R=1000$.

\section{Metodika stanovení rozsahu výběru}
Součástí analýzy je teoretický výpočet potřebného rozsahu výběru $n$ pro dosažení požadované přesnosti $\Delta$. Vycházíme z požadavku, aby poloviční šířka intervalu spolehlivosti nepřesáhla $\Delta$.
Z asymptotického rozptylu (3.3) plyne:
\begin{equation}
n \ge \frac{u_{1-\alpha/2}^2 p(1-p)}{\Delta^2 [f(Q_p)]^2}
\end{equation}
Pro log-normální rozdělení je hustota dána vzorcem:
\begin{equation}
f(x) = \frac{1}{x\sigma\sqrt{2\pi}} e^{-\frac{(\ln x - \mu)^2}{2\sigma^2}}
\end{equation}
Tento vzorec ukazuje, že s rostoucím $x$ (v chvostu u $\sigma=1.5$) hustota $f(x)$ rychle klesá k nule. Protože $f(Q_p)$ je ve jmenovateli vzorce pro $n$, potřebný rozsah výběru roste s druhou mocninou převrácené hodnoty hustoty. To má dramatické důsledky pro plánování šetření zaměřených na extrémní kvantily.
